\chapter{犯罪主体}

% Source PDF physical page 36; see audit/source_coverage.json.
\begin{EditionOpeningPanel}
考查重点是：定罪身份（国家工作人员），单位犯罪。
\end{EditionOpeningPanel}

\section{自然人}

[问题]甲是公立医院的副院长，周一，与医疗器材商洽谈采购事宜，收受对方财物；周二，坐门诊，开处方，收受医药代表的钱财，答应多开其药品。甲构成何罪？（副院长收钱案）

自然人要成为犯罪主体（行为主体），只要求具有侵害法益的能力（违法能力）即可，不要求具有责任能力，因为责任能力是可谴责性的条件，与制造违法事实没有关系。

例如，精神病患者、13岁的人，均能制造违法事实（例如杀人），均是制造违法事实意义上的犯罪主体；至于责任能力、责任年龄，则是后面的排除责任事由。关于自然人，主要问题是特殊身份。特殊身份，是指行为人在身份上的特殊资格或特殊地位。

\subsection*{一、真正身份犯}

这是指行为人只有具备某种特殊身份，才能构成犯罪。这种特殊身份也称为定罪身份或构成身份。例如，刑讯逼供罪的主体必须是司法工作人员，贪污罪的主体必须是国家工作人员。

[注意1]定罪身份必须在开始犯罪时就具有。如果是在犯罪过程中形成的身份，则不属于定罪身份。例如，组织、领导、参加黑社会性质组织罪的组织者，不属于定罪身份，不是真正身份犯。因为这种身份是在犯罪过程中形成的。又如，犯罪集团中的首要分子，也不是定罪身份。（2020年试题）

[注意2]定罪身份只是针对实行犯所要求的。要构成共犯（帮助犯、教唆犯），不要求有定罪身份。例如，官员的妻子是普通公民，不可能构成贪污罪的实行犯，但可以构成贪污罪的教唆犯、帮助犯。

\subsection*{二、不真正身份犯}

这是指行为人具有某种特殊身份，不影响犯罪的成立，但是影响量刑。这种特殊身份也称为量刑身份或加减身份。例如，诬告陷害罪的主体是一般主体，但是国家机关工作人

% Source PDF physical page 37; see audit/source_coverage.json.
员犯该罪就要从重处罚。诬告陷害罪就是不真正身份犯，国家机关工作人员的身份就是量刑身份。\footnote{国家机关工作人员的身份，在报复陷害罪（第254条）中是定罪身份，而在诬告陷害罪（第243条）中是量刑身份。注意：打击报复证人罪（第308条）不是身份犯，不要求行为人有身份。}

\subsection*{三、“国家工作人员”的认定}

\InsertDiagram{figures/manual_table_p037_1.tex}

错误标准是，根据有无正式编制来认定，往往认为临时工不是国家工作人员。然而，根据刑法第93条规定，正确标准是：是否从事公务。\footnote{根据刑法第93条的规定和立法解释，国家工作人员包括四类人：（1）国家机关中从事公务的人员；（2）国有公司、企业、事业单位、人民团体中从事公务的人员。人民团体不包括社会团体。常见的人民团体有妇联、共青团等；（3）国家机关、国有公司、企业、事业单位委派到非国有公司、企业、事业单位、社会团体从事公务的人员；（4）其他依照法律从事公务的人员。根据立法解释，村民委员会等村基层组织人员协助政府从事下列行政管理工作（公务）时，属于上述（4）中的国家工作人员：①救灾、抢险、防汛、优抚、扶贫、移民、救济款物的管理；②社会捐助公益事业款物的管理；③国有土地的经营和管理；④土地征收、征用补偿费用的管理（注意：不是指土地承包费）；③代征、代缴税款；6有关计划生育、户籍、征兵工作。}

\CJKunderwave{判断公务的标准}：第一，事务具有公共管理性。这是指事务关系到多数人或不特定人的利益。仅与个别人或少数人相关的事务，不是公务。第二，事务具有行政职责性。这是指事务属于行政职务，并承担行政责任。例如，公立大学财务处会计属于国家工作人员，大学教授属于非国家工作人员。又如，村干部协助政府从事公务时，属于国家工作人员。

例如，本节开头的“副院长收钱案”中，周一，甲从事的是公务，是国家工作人员，构成受贿罪；周二，甲从事的是技术劳动，是非国家工作人员，构成非国家工作人员受贿罪。（2016年第62题）

\section{单位犯罪}

[问题］狗蛋问厂长：“我想干掉情敌狗剩，你的化工厂有氰化钾，借给我点。”该厂领导层决议，答应帮这个忙。狗蛋拿到氰化钾，杀死了狗剩。对该厂的行为如何处理？（化工厂协助杀人案）

\subsection*{一、分类}

1.纯正的单位犯罪。这是指只能由单位构成而不能由自然人构成的犯罪。例如，单位行贿罪、单位受贿罪。

2.不纯正的单位犯罪。这是指既可由单位构成也可由自然人构成的犯罪。例如，生产、销售伪劣产品罪。\footnote{不纯正的单位犯罪主要集中在“破坏经济秩序犯罪”一讲。}

［提示］纯正的自然人犯罪，是指只能由自然人构成而不能由单位构成的犯罪。例如，

% Source PDF physical page 38; see audit/source_coverage.json.
行贿罪、受贿罪。又如，金融诈骗罪的8个罪名中，只有3个罪名是纯正的自然人犯罪：\CJKunderwave{贷款}诈骗罪、\CJKunderwave{信用}卡诈骗罪、\CJKunderwave{有价}证券诈骗罪。

\subsection*{二、成立条件}

\subsubsection*{（一）主体条件（法人资格问题）}

1.单位犯罪是否要求单位必须具有法人资格？

[结论］国有单位，不要求法人资格；私营单位，要求法人资格。例如，私营公司有法人资格，可以构成单位犯罪；合伙没有法人资格，不能构成单位犯罪。

2.单位的分支机构或内设机构，没有独立法人资格，要成为单位犯罪的主体，需要额外满足两个条件：（1）以自己名义犯罪；（2）违法所得归该机构所有。\footnote{这两个额外条件表明，没有法人资格的国有单位并非一律都能成为单位犯罪的主体。像这种分支机构、内设机构要成为单位犯罪主体，还是有额外要求的。}（2006年第5题，2020年试题）

例如，分公司不具有独立法人资格，属于公司的分支机构，但具备上述两个条件，也能成为单位犯罪的主体。子公司有独立法人资格，能成为单位犯罪的主体。（2019年客观题）

\subsubsection*{（二）主观条件}

1.成立单位犯罪，要求具有单位的意志。单位犯罪的犯罪意志不是内部某个成员的意志，而是单位的整体意志。单位整体意志的产生方式有：第一，由单位决策机构形成（如董事会决议）；第二，单位领导依据职权作出决策。

2.关于主观罪过形式，单位犯罪既可以是故意犯罪，也可以是过失犯罪。单位过失犯罪，例如，工程重大安全事故罪（第137条）、出具证明文件重大失实罪（第229条第3款）。

3.单位犯罪的客观样态与自然人犯罪的客观样态没有区别，都是由自然人去实施。

4.揭开单位的面纱。根据司法解释，两种单位的犯罪，按照自然人犯罪论处：（1)成立单位时\CJKunderwave{主要目的}就是犯罪；（2）成立单位后\CJKunderwave{主要活动}就是犯罪。亦即，此时“单位”属于个人犯罪的“马甲”，应当揭开“马甲”，直接定个人犯罪。此处原理，类似于公司法中的“揭开公司的面纱”。（2006年第5题）

\subsection*{三、单位犯罪与个人犯罪的关系}

\subsubsection*{（一）如何区分“单位犯罪”与“单位内部成员的个人犯罪”？}

1.主观条件：单位犯罪体现单位意志，内部成员个人犯罪体现个人意志。未经单位授权或未依据职权作出的决策，不能体现单位意志，只能体现个人意志。

2.客观条件：单位犯罪，是为单位（全体成员）谋取非法利益。内部成员个人犯罪，是为特定成员个人谋取非法利益。

% Source PDF physical page 39; see audit/source_coverage.json.
\InsertDiagram{figures/manual_p039_1.tex}

（1）模型1，单位意志+本单位利益。

例如，甲公司董事会决议，向官员乙行贿，以便为公司谋取非法利益。这是单位行贿罪。

（2）模型2，单位意志+成员个人利益。

例1，甲公司董事会决议，向官员乙行贿，以便为甲公司董事长的孩子办理升学违规手续。这属于为单位内部特定成员谋取非法利益，不是为单位谋取非法利益，不构成单位犯罪（单位行贿罪），而构成个人犯罪（行贿罪），行为主体是公司的直接责任人员。

例2（2008年主观题），某大型国有公司的领导层私分了公司一笔应收账款100万元。该犯罪虽然体现了单位意志，但未为单位谋取利益，而是为单位内部特定成员谋取利益，构成成员的个人犯罪，也即贪污罪，而不构成单位犯罪（私分国有资产罪）。

（3）模型3，成员个人意志+本单位利益。

例1，甲公司的副总王某向董事会提议，向官员乙行贿，才能为公司谋取到非法利益，但董事会不同意。王某私下用自已的钱向乙行贿，为公司谋取非法利益。甲公司董事会对此并不知情，未体现单位意志，因此这是成员的个人犯罪（行贿罪）。由此也可看出，个人犯罪并不要求必须为本人谋取非法利益，也可以为本人以外第三方（包括本单位、外单位、第三人）谋取非法利益。这如同，盗窃罪的非法占有目的，既包括为自己非法占有，也包括为第三方非法占有。

例2（2015年第54题），某公司的经理赵某，为了单位利益，实施拒不执行判决行为，但该起犯罪活动未体现单位意志。由于未体现单位意志，因此赵某构成个人犯罪。

（4）模型4，成员个人意志+个人利益。

例1，甲公司经理李某盗用公司名义，实施走私犯罪。李某构成个人犯罪。

例2，国有公司的副总王某未经单位授权，向官员行贿，以便为自已谋取非法利益。王某构成个人犯罪（行贿罪）。\footnote{需注意“为本单位谋取非法利益”条件的应用场景。这个条件仅在区分单位犯罪与单位成员个人犯罪时有意义。如果区分单位犯罪与单位之外的个人犯罪，则不需要考虑这个条件。此时，成立单位犯罪，不要求为本单位谋取非法利益，可以为其他单位或单位成员外的个人谋取非法利益。这就如同普通的个人犯罪一样，并不要求必须为本人谋取非法利益，也可以为本人以外第三方谋取非法利益。例1，赵某向官员行贿，目的是为朋友李某谋取非法利益。赵某构成行贿罪。例2，甲公司董事会决议，向官员行贿，以便为乙公司或乙公司的领导谋取非法利益。甲公司成立单位行贿罪。}

\subsubsection*{（二）单位与个人的共同犯罪问题}

1.单位犯罪是单位本身的犯罪，不是各个成员之间的共同犯罪，也不是单位与成员之间的共同犯罪。（2020年试题）

2.甲单位与乙单位之间可以构成共同犯罪。单位与单位外自然人之间可以构成共同犯罪。

% Source PDF physical page 40; see audit/source_coverage.json.
\subsubsection*{（三）单位实施纯正的个人犯罪}

问题：单位实施只能由个人构成的犯罪（纯正的个人犯罪），如何处理？结论：定单位的直接责任人的个人犯罪。\footnote{对此，法律依据是2014年4月24日全国人民代表大会常务委员会《关于\ensuremath{<}中华人民共和国刑法\ensuremath{>}第三十条的解释》规定：“公司、企业、事业单位、机关、团体等单位实施刑法规定的危害社会的行为，刑法分则和其他法律未规定追究单位的刑事责任的，对组织、策划、实施该危害社会行为的人依法追究刑事责任。”}

例1，某个单位领导集体决议实施盗电，为单位谋取非法利益。虽然刑法规定单位不能成为盗窃罪主体，但是可以直接追究直接责任人的自然人犯罪，也即盗窃罪。自然人犯罪并不要求必须为自然人本人谋取非法利益，也可以为第三人（某个单位或个人）谋取非法利益。

例2，刑法规定贷款诈骗罪只能由自然人构成。如果单位实施贷款诈骗，虽然不能认定单位构成贷款诈骗罪，但可以认定单位的直接责任人构成贷款诈骗罪。\footnote{提示：以前有司法解释认为，单位实施贷款诈骗的，只能以单位构成合同诈骗罪来处理。这是错误的。2014年立法解释已经纠正了该错误。}这是三段论推导的结论，大前提（贷款诈骗罪只能由自然人构成），小前提（甲公司实施贷款诈骗，由直接责任人具体实施），结论：直接责任人构成贷款诈骗罪。（2011年第12题）

例3，刑法规定伪造货币罪只能由自然人构成。某单位伪造货币，直接追究相关自然人的伪造货币罪。（2015年第54题）

[引申]单位帮助自然人实施纯正的自然人犯罪，能否对单位以帮助犯论处？不能，只能追究其直接责任人的帮助犯责任。

例1，星星公司领导层得知赵某欲实施贷款诈骗，便集体决定帮助赵某，为赵某提供了虚假担保。赵某诈骗银行贷款得逞。赵某构成贷款诈骗罪的实行犯，星星公司直接责任人构成贷款诈骗罪的帮助犯，但不能认定星星公司构成贷款诈骗罪的帮助犯。

例2，在本节开头“化工厂协助杀人案”中，化工厂的直接责任人构成狗蛋的故意杀人罪的帮助犯。

\subsubsection*{（四）单位实施不纯正的单位犯罪}

不纯正的单位犯罪，是指既可由单位构成也可由自然人构成的犯罪。注意：就同一犯罪而言，虽然行为主体不同，但是定罪标准、既遂标准应当相同。不过，有些司法解释并未遵守这一点，往往将单位犯罪的定罪标准（如定罪数额）定得高一些，将自然人犯罪的定罪标准（如定罪数额）定得低一些。只能将这些做法视为特殊规定。

问题：单位实施单位和个人都能构成的犯罪（不纯正的单位犯罪），单位未达到单位犯罪的定罪数额标准，能否认定直接责任人的个人犯罪？结论：不能。

例如，走私普通货物、物品罪既可以由单位构成，也可以由个人构成。司法解释规定，单位的定罪数额标准是20万元，个人的定罪数额标准是10万元。甲单位走私数额是15万元，不构成单位犯罪，对此也不能定单位的直接责任人的走私普通货物、物品罪。

理由是，既然刑法规定一个罪名既可以由单位构成，也可以由个人构成，并且为二者规定了不同的定罪数额标准，那么二者便是独立关系，并水不犯河水，不能从单位向个人“串线”。

上述单位盗电，定直接责任人的盗窃罪，这是就纯正的个人犯罪而言的。刑法规定盗窃罪只能由个人构成，现在单位实施盗窃，没法追究单位，只能“串线”到个人犯罪。

% Source PDF physical page 41; see audit/source_coverage.json.
记结论：单位实施纯正个人犯罪，允许“串线”到个人犯罪；单位实施不纯正单位犯罪，不允许“串线”到个人犯罪。（2019年试题）

\subsection*{四、处罚}

\subsubsection*{（一）处罚规则}

1.原则双罚，既罚单位，也罚个人（主管人员和直接责任人）。

（1）对单位本身只能判处罚金，不能判处没收财产。\footnote{对单位不判处没收财产，主要是考虑到没收财产是指没收合法的没有用于犯罪的财产，如果没收单位的财产如办公设备、生产资料等，会影响单位的正常经营。单位犯了罪，并不意味着该单位就必须被撤销。单位犯了罪，不意味着单位变成了犯罪集团。}（2010年第53题）

（2）所处罚的个人，必须是参与犯罪的人。不能因为单位构成犯罪，便处罚没有参与犯罪的主管人员。

（3）所处罚的个人，只能是两类人（主管人员和直接责任人）。\CJKunderwave{其他人即使参与了单}\CJKunderwave{位犯罪，也不得处罚}。

（4）处罚个人时，数个个人之间可以构成共同犯罪。但该共同犯罪是处在单位犯罪的框架下的，是为了明确各个人的责任大小、量刑轻重。

2.例外单罚，只处罚个人（主管人员和直接责任人员），不处罚单位。这是因为，如果处罚单位，会损害无辜的普通职员的利益，有株连无辜之嫌。\footnote{不需要特别去记例外单罚的具体情形和罪名，因为考试很少考到这个层级。法考刑法很少考纯粹死记硬背的东西。这一点和法考刑诉法的考查理念有所不同。}（2010年第53题）

\subsubsection*{（二）“单位没了”}

\InsertDiagram{figures/manual_table_p041_1.tex}

\OriginalPanel{典型真题}

关于单位犯罪，下列哪些选项是正确的？（2015年·卷二·54题）\footnote{[答案]A项，该说法是针对正常原理而言的，所以该说法正确。该说法没有考虑某些司法解释的特别规定。C项，刑法第313条规定的拒不执行判决罪，在2015年11月1日生效的《刑法修正案（九）》之前，只能由自然人构成，生效后，则可以由单位构成。本题为2015年的考题，所以说法错误。再者，即使本案发生在2015年11月1日后，也不构成单位犯罪，因为赵某的行为虽为单位谋利，但未体现单位的整体意志。本题答案：AD。}

A.就同一犯罪而言，单位犯罪与自然人犯罪的既遂标准完全相同

B.《刑法》第一百七十条未将单位规定为伪造货币罪的主体，故单位伪造货币的，相关自然人不构成犯罪

C.经理赵某为维护公司利益，召集单位员工殴打法院执行工作人员，拒不执行生效判决的，成立单位犯罪

D.公司被吊销营业执照后，发现其曾销售伪劣产品20万元。对此，应追究相关自然人销售伪劣产品罪的刑事责任

